\documentclass[12pt]{article}
\usepackage[utf8]{inputenc}
\usepackage[T1]{fontenc}
\usepackage[margin=1in]{geometry}
\usepackage{setspace}
\usepackage{xcolor}
\usepackage[colorlinks=true, urlcolor=blue!60!black, linkcolor=blue!60!black]{hyperref}
\setlength{\parindent}{0em}
\setlength{\parskip}{0.45em}
\renewcommand{\baselinestretch}{1.0}
\begin{document}
\begin{center}
{\Large \textit{Loving v.\ Virginia}, 388 U.S.\ 1 (1967)\par}
\vspace{0.3em}
{\normalsize Case Brief --- One-Page Edition\par}
\end{center}

\textbf{Memory Jogger.} Virginia's interracial-marriage ban violated equal protection and due process (\href{https://supreme.justia.com/cases/federal/us/388/1/}{Justia}).

\textbf{Facts.} Mildred Jeter, a Black woman, and Richard Loving, a white man, married in Washington, D.C., in 1958 and returned to Virginia, which made interracial marriage a felony. Police arrested them in a night raid.

\textbf{Procedural History.} They pleaded guilty and were sentenced to a year in jail, suspended if they left Virginia; Virginia's high court upheld the statutes.

\textbf{Votes.} Reversed 9--0: Warren for the Court; Stewart concurred in the judgment.

\textbf{Holding.} Racial classifications in marriage, even if ``equally applied,'' face ``the most rigid scrutiny'' and fail; and the freedom to marry is a fundamental liberty race cannot be used to deny.

\textbf{Reasoning.} Equal application does not immunize a racial classification, and the framers' debates were ``at best, inconclusive.'' Banning only marriages involving white persons showed the laws were ``designed to maintain White Supremacy.'' Marriage is one of the ``basic civil rights of man.'' Justice Stewart would hold race-based criminal liability invalid per se.

\textbf{Examples.} \emph{Same-Side:} a ban on marriage between people of different national origins, applied equally to all groups. \emph{Opposite-Side:} licensing, minimum-age, and consanguinity rules. \emph{Fence-Sitter:} a bar on marriage for parents behind on child support, later struck down in \textit{Zablocki v.\ Redhail}.

\textbf{Critique.} Many framers expected such bans to survive, making \textit{Loving} a test case for originalism.

\textbf{Key Quotation.} ``The freedom to marry has long been recognized as one of the vital personal rights essential to the orderly pursuit of happiness by free men.'' --- Chief Justice Warren, 388 U.S.\ at 12.

\end{document}
